% ex: ts=2 sw=2 sts=2 et filetype=tex
% SPDX-License-Identifier: CC-BY-SA-4.0
\documentclass[aspectratio=169]{beamer}

\usetheme{Boadilla}

\setbeamertemplate{navigation symbols}{} % To remove the navigation symbols from the bottom of all slides

\graphicspath{{../../../slides/python/img/}{../styles/iteso/}}

% use the command "usepackage" instead of "usecolortheme" to use the path to styles
\usepackage{../styles/iteso/beamercolorthemeiteso}

\usepackage[utf8]{inputenc}
\usepackage[T1]{fontenc} %% https://tex.stackexchange.com/questions/664/why-should-i-use-usepackaget1fontenc
\usepackage{graphicx} % Allows including images
\usepackage{listings}

\title{Entrada y Salida de datos de usuario} % The short title appears at the bottom of every slide,
                                               % the full title is only on the title page
\author{Miguel Bernal Marin}
\institute[ITESO]
{
 ITESO, Universidad \\
 Jesuita de Guadalajara \\
\medskip
\textit{miguel.bernal@iteso.mx}
}
\date{
Otoño 2024
} % Date, can be changed to a custom date

\newcommand{\pausa}{\pause} % Para usar una pausa en las presentaciones
%\newcommand{\pausa}{}      % Para que NO salgan las pausas

%New colors defined below
\definecolor{codeBakground}{rgb}{0.95,0.95,0.92}
\definecolor{codeComment}{rgb}{0.0, 0.0, 1.0}
\definecolor{codeKeyword}{rgb}{0.647, 0.165, 0.165}
\definecolor{codeKeyword2}{rgb}{0.0, 0.545,0.545}
\definecolor{codeNumbers}{rgb}{0.5,0.5,0.5}
\definecolor{codeString}{rgb}{1.0, 0.0, 1.0}
\definecolor{light-gray}{gray}{0.90}

%Code listing style named "mystyle"
\lstdefinestyle{mystyle}{
  backgroundcolor=\color{codeBakground},
  commentstyle=\color{codeComment},
  keywordstyle=\color{codeKeyword},
  keywordstyle={[2]\color{codeKeyword2}}, % Built-in
  numberstyle=\tiny\color{codeNumbers},
  stringstyle=\color{codeString},
  basicstyle=\ttfamily\footnotesize,
  breakatwhitespace=false,
  breaklines=true,
  captionpos=b,
  keepspaces=true,
  numbers=left,
  numbersep=5pt,
  showspaces=false,
  showstringspaces=false,
  showtabs=false,
  upquote=true,                      % requires textcomp
  tabsize=2,
  % Español
  literate=%
  {á}{{\'a}}1 {é}{{\'e}}1 {í}{{\'i}}1 {ó}{{\'o}}1 {ú}{{\'u}}1
  {Á}{{\'A}}1 {É}{{\'E}}1 {Í}{{\'I}}1 {Ó}{{\'O}}1 {Ú}{{\'U}}1
  {¡}{{!`}}1  {¿}{{?`}}1
  {ñ}{{\~n}}1 {Ñ}{{\~N}}1
}

%"mystyle" code listing set
\lstset{style=mystyle}

%------------------------------------------------------------
%The next block of commands puts the table of contents at the 
%beginning of each section and highlights the current section:
\AtBeginSection[]
{
  \begin{frame}
    \frametitle{Contenido}
    \tableofcontents[currentsection]
  \end{frame}
}
%------------------------------------------------------------

\begin{document}
{ % Portada de la universidad
  \setbeamertemplate{footline}{}
  \usebackgroundtemplate{\includegraphics[width=\paperwidth]{fondo-portada.png}}
  \begin{frame}
  \end{frame}
}

\begin{frame}
    \titlepage
\end{frame}

\usebackgroundtemplate{\includegraphics[width=\paperwidth]{fondo.png}}

\begin{frame}
    \frametitle{Contenido}
    \tableofcontents
\end{frame}

% ex: ts=4 sw=4 sts=4 et filetype=tex
% SPDX-License-Identifier: CC-BY-SA-4.0
\begin{frame}
    \frametitle{Contenido}
    \tableofcontents
\end{frame}

\section{Algoritmo y sus características}

\begin{frame}[c]{Título}
    \begin{center}
        Texto
    \end{center}
\end{frame}

\section{Lenguajes de programación y su clasifiación}

\begin{frame}[c]{Título}
    \begin{center}
        Texto
    \end{center}
\end{frame}


{ % Portada de la universidad
  \setbeamertemplate{footline}{}
  \usebackgroundtemplate{\includegraphics[width=\paperwidth]{fondo-portada.png}}
  \begin{frame}
  \end{frame}
}
\end{document}
